% =====================================================================
%  fm.tex ------ 前置：まえがき（序）+ 理論の概要と展望（理论的概要与展望）
%  原书 pp.v--xviii（PDF 5--18）；Lead 亲译，作为全书文风样板
% =====================================================================

\chapter*{序}
\addcontentsline{toc}{chapter}{序}
\markboth{序}{序}

% [原文 p.v]
物理学与几何学历来相互给予强烈的影响而共同发展。1980 年代以来，两者关系中最显著的一点，是量子场论在几何学上的应用。本书所处理的，是以 1980 年代 Jones 多项式的发现、共形场理论、Witten 用 Chern--Simons 规范理论构造三维流形的不变量等为契机而发展起来的，二维量子场论与三维拓扑学之间的关系。

量子场论的本质困难，在于其自由度为无限大。量子场论中发展起来的处理无穷的种种手法，也已深深渗透到几何学之中。这里，我们把注意力集中在「以流形上场的配分函数来构造流形的拓扑不变量」这一点上，叙述它向拓扑学的应用。

场论中有几种被物理学家视为彼此等价的表述，例如 Feynman 路径积分的量子化与正则量子化。把几何学所处理的对象翻译成场论的概念，有时会发现新的联系、带来意想不到的展开。本书的主题------共形场理论与 Chern--Simons 理论的关系------也正是从 Chern--Simons 泛函出发、用两种方法写出态空间之后才被发现的。这样的手法为几何学提供了新的视角，带来种种启示。不过需要记住：物理意义上的等价性，在几何学中未必是被证明了的事实。为了证明它，往往需要个别地开发手法，而这又向几何学一方提出了重要的问题。

以下概观本书的内容。从古典力学开始，关于几何量子化等背景性的基础概念，在「理论的概要与展望」中叙述。第 1 章以环路群的几何学为基础处理共形场理论。

% [原文 p.vi]
首先解说环路群及其上的复线丛、仿射 Lie 代数、Virasoro 代数等，然后对共形场理论中的重要对象------共形块空间、KZ 方程等加以表述。KZ 方程通常由关于场算子的算子积展开导出，这里则强调几何的视角来叙述。第 2 章解说由共形场理论的单值（holonomy）得到的辫群表示，并用它构造链环的 Jones 多项式与三维流形的 Witten 不变量。此外还描述曲面映射类群对共形块空间的作用，以及 Atiyah 意义下的拓扑场论。2.5 节触及 Chern--Simons 理论与曲面平坦联络的模空间的几何量子化之间的关系。第 3 章处理 Chern--Simons 微扰理论与 Vassiliev 不变量。先讲述纽结的 Vassiliev 不变量，说明从 Chen 的迭代积分到 Kontsevich 的 Vassiliev 不变量积分表示这一脉络。最后描述由 Chern--Simons 微扰理论导出的三维流形拓扑不变量。

本书所处理的内容，从场论与拓扑学的关系来看不过是极小的一部分；若读者能通过这里处理的题材感受到理论所蕴含的丰富性，对笔者而言便是望外之喜。

\medskip
\hspace*{2em}1998 年 7 月

\hspace*{2em}河野俊丈

\bigskip
\noindent\textbf{追记}\quad 本书是将作为岩波讲座《现代数学的展开》之一分册刊行的《场论与拓扑学》单行本化的产物。此次单行本化之际，作了一部分订正。

\medskip
\hspace*{2em}2008 年 3 月

\chapter*{理论的概要与展望}
\addcontentsline{toc}{chapter}{理论的概要与展望}
\markboth{理论的概要与展望}{理论的概要与展望}

% [原文 p.vii]
本书的目的是介绍受场论中各种手法启发而发展起来的拓扑学研究的新进展。为此，先从古典力学开始，叙述理论的背景与基础概念。

\para{古典力学中的 Lagrange 形式与 Hamilton 形式}
设 $N$ 为 $n$ 维微分流形。设在 $N$ 的切丛 $TN$ 上给定了依赖于时间参数 $t$、称为 Lagrangian 的函数 $\mathcal{L}(t,x,\xi)$。这里 $x=(x_1,\dots,x_n)$ 是 $N$ 的局部坐标，$\xi=(\xi_1,\dots,\xi_n)$ 是切丛坐标。设 $\gamma\colon[a,b]\to N$ 为 $N$ 的光滑曲线，称
\[
  S=\int_a^b \mathcal{L}\bigl(t,\gamma(t),\gamma'(t)\bigr)\,dt
\]
为作用量积分。固定端点 $\gamma(a),\gamma(b)$ 而变动曲线时，使作用量积分取临界值的曲线满足 Euler--Lagrange 方程
\[
  \frac{d}{dt}\frac{\partial \mathcal{L}}{\partial \xi_j}
  =\frac{\partial \mathcal{L}}{\partial x_j},\qquad j=1,\dots,n .
\]
特别地，对 $n$ 维 Euclid 空间中质量为 $m$ 的质点的运动，取势为 $V$、令
$\mathcal{L}=\frac{1}{2}m\sum_{j=1}^{n}\xi_j^{2}-V$，则 Euler--Lagrange 方程无非就是 Newton 运动方程。像这样，以 Lagrangian 为出发点、把运动方程描述为作用量积分的变分问题的解，这种古典力学的表述方式称为 Lagrange 形式。

其次说明 Hamilton 形式。一般地，设 $M$ 为 $2n$ 维流形。$M$ 上满足 $d\omega=0$ 且非退化的 2 次微分形式 $\omega$ 称为\textbf{辛形式}（symplectic form）。这里所谓的 2 次微分形式

% [原文 p.viii]
非退化，是指 $M$ 上处处有 $\omega^{n}\neq 0$。带有辛形式的流形称为\textbf{辛流形}（symplectic manifold）。

一旦有了辛形式，就可以在向量场与 1 次微分形式之间建立一一对应。设 $X$ 为 $M$ 上的向量场，定义 1 次微分形式 $\varphi=i(X)\omega$ 为
\[
  \varphi(Y)=\omega(X,Y),
\]
这里 $Y$ 是 $M$ 上的向量场。由辛形式的非退化性，上述对应
\[
  X\longmapsto i(X)\omega
\]
给出向量场与 1 次微分形式之间的一一对应。对 $M$ 上的光滑函数 $f$，定义向量场 $X_f$ 为
\[
  i(X_f)\omega=df,
\]
称它为对应于 $f$ 的 \textbf{Hamilton 向量场}（Hamiltonian vector field）。

对 $M$ 上的函数 $f,g$，定义 $\{f,g\}$ 为
\[
  \{f,g\}=-\omega(X_f,X_g),
\]
称它为 $f$ 与 $g$ 的 \textbf{Poisson 括号}（Poisson bracket）。Poisson 括号关于 $f,g$ 是线性的、反对称的。由定义，
\[
  \omega(X_f,X_g)=i(X_f)\omega(X_g)=df(X_g)=X_g f,
\]
它也写作 $-X_f g$。又，对 Hamilton 向量场的 Lie 括号，有 $[X_f,X_g]=X_{\{f,g\}}$ 成立。

设在 $M$ 上给定了称为 Hamiltonian 的函数 $H$，考虑 $H$ 对应的 Hamilton 向量场 $X_H$。我们来描述 $M$ 上函数 $f$ 的值沿 Hamilton 向量场的积分曲线如何变化。以 $t$ 为积分曲线的参数，则
\[
  \frac{df}{dt}=X_H f,
\]
因此

% [原文 p.ix]
\[
  \frac{df}{dt}=\{H,f\}
\]
成立。特别地，Hamiltonian $H$ 在上述积分曲线上取常值。

\para{复线丛与量子化}
设 $E$ 为微分流形 $M$ 上的 $C^{\infty}$ 级复向量丛。$M$ 的余切丛 $T^{*}M$ 的复化记作 $T^{*}M_{\C}$。$E$ 与 $T^{*}M_{\C}\otimes E$ 的 $C^{\infty}$ 级截面全体所成的空间分别记作 $\Gamma(E)$、$\Gamma(T^{*}M_{\C}\otimes E)$。$E$ 上的\textbf{联络}（connection）是指 $C$ 线性映射
\[
  \nabla\colon \Gamma(E)\longrightarrow \Gamma(T^{*}M_{\C}\otimes E),
\]
它使得对 $M$ 上的 $C^{\infty}$ 级函数 $f$ 与 $E$ 的 $C^{\infty}$ 级截面 $s$，Leibniz 法则
\[
  \nabla(fs)=df\otimes s+f\nabla(s)
\]
成立。对向量场 $X\in\Gamma(TM_{\C})$，由
\[
  (\nabla_X s)(x)=(\nabla s)(X_x)
\]
确定的 $\nabla_X\colon \Gamma(E)\to\Gamma(E)$，称为沿 $X$ 的\textbf{协变微分}（covariant derivative）。关于向量丛与联络的一般理论，请参考文献 [6]。

以下只说明秩为 $1$ 的复向量丛，即\textbf{复线丛}（complex line bundle）。设 $M$ 为微分流形，$L$ 为其上的 $C^{\infty}$ 级复线丛，并设 $L$ 上配有 Hermite 度量。存在 $M$ 的开覆盖 $M=\bigcup_{j}U_j$，使得在开集 $U_j$ 上 $L$ 被平凡化为 $U_j\times\C$。设 $\nabla$ 为 $L$ 的联络，在开集 $U_j$ 上写成
\[
  \nabla=d-2\pi\sqrt{-1}\,\alpha_j .
\]
这里 $\alpha_j$ 是 $U_j$ 上定义的局部 1 次微分形式，而 $d\alpha_j$ 确定 $M$ 上的整体 2 次微分形式；把它记作 $c_1(\nabla)$，称为 $\nabla$ 的\textbf{第一 Chern 形式}（first Chern form）。de Rham 上同调类 $[c_1(\nabla)]\in H^{2}(M,\R)$ 不依赖于联络 $\nabla$ 的取法，这就是 $L$ 的第一 Chern 类。\textbf{第一 Chern 类}（first Chern class）包含在自然映射

% [原文 p.x]
$i\colon H^{2}(M,\Z)\to H^{2}(M,\R)$ 的像中。反之，取 $M$ 上满足 $d\omega=0$ 且确定 $H^{2}(M,\Z)$ 中元素的 2 次微分形式 $\omega$，则可以构造 $L$ 的 Hermite 联络 $\nabla$，使 $c_1(\nabla)=\omega$。第 1 章将对 $M$ 单连通的情形实际构造复线丛 $L$。

在古典力学中，表示系统时间发展的运动方程，用 Hamiltonian $H$ 与 Poisson 括号写成 $\frac{df}{dt}=\{H,f\}$。Dirac 的量子力学框架则是：对函数 $f$ 对应一个线性算子 $\hat f$，要求这一对应是线性的，并且当 $\{f_1,f_2\}=f_3$ 时，
\[
  \hat f_1\hat f_2-\hat f_2\hat f_1=-\sqrt{-1}\,\hbar\,\hat f_3
\]
成立。这里令 $\sqrt{-1}\hbar\,\tilde f/\hbar=\hat f$。算子 $\hat f_1,\hat f_2$ 不可交换，用
\[
  [\hat f_1,\hat f_2]=\hat f_1\hat f_2-\hat f_2\hat f_1
\]
定义算子的 Lie 括号。量子力学的基本方程是
\[
  \frac{d\hat f}{dt}=[\hat H,\hat f],
\]
它被称为 \textbf{Heisenberg 方程}。

把古典力学系统表示为辛流形 $(M,\omega)$，把 $M$ 上光滑函数全体 $C^{\infty}(M)$ 关于 Poisson 积看作 Lie 代数。把这个 Lie 代数表示在某个 Hilbert 空间上------也就是说，对每个函数 $f$ 对应一个 Hilbert 空间的线性算子，并保持 Lie 代数结构------这就是\textbf{量子化}（quantization）的基本想法。我们考虑用 $M$ 上的复线丛来实行这一纲领。如上所述，在「$\omega$ 确定整系数上同调 $H^{2}(M,\Z)$ 中的元素」这一假设之下，可以构造 $M$ 上的复线丛 $L$ 及其上的联络 $\nabla$，使对应的第一 Chern 形式等于 $\omega$。

把复线丛 $L$ 的光滑截面全体记作 $\Gamma(M,L)$，对 $s_1,s_2\in\Gamma(M,L)$，令

% [原文 p.xi]
\[
  \langle s_1,s_2\rangle=\int_M \bigl(s_1(x),s_2(x)\bigr)\,\frac{\omega^{n}}{n!},
\]
由此引入 $L^{2}$ 内积。这里上式右端的括号表示复线丛 $L$ 的 Hermite 度量。关于这个内积，把 $L$ 的 $L^{2}$ 截面全体的空间记作 $\mathcal{H}$，则 $\mathcal{H}$ 具有 Hilbert 空间的结构。以 $\nabla_{X_f}$ 表示沿 Hamilton 向量场 $X_f$ 的协变微分。对 $M$ 上的光滑函数 $f$，定义 $\mathcal{H}$ 的线性算子 $\hat f$ 为
\[
  \hat f s=\nabla_{X_f}s-2\pi\sqrt{-1}\,f s,\qquad s\in\mathcal{H}.
\]
这里，利用
\[
  \nabla_X\nabla_Y-\nabla_Y\nabla_X=\nabla_{[X,Y]}-2\pi\sqrt{-1}\,\omega(X,Y)
\]
与 $\omega(X_f,X_g)=-\{f,g\}$，由上述 $\hat f$ 的构造可知 $[\hat f,\hat g]=\widehat{\{f,g\}}$ 成立。也就是说，可以确认：关于 Poisson 括号的 Lie 代数 $C^{\infty}(M)$ 在 $\mathcal{H}$ 上的表示被得到。

取微分流形 $N$，令余切丛 $T^{*}N$ 为 $M$。用 $N$ 的局部坐标 $q_1,\dots,q_n$ 与纤维方向的坐标 $p_1,\dots,p_n$，$M$ 成为带有辛形式
\[
  \omega=\sum_{j=1}^{n} dp_j\wedge dq_j
\]
的辛流形。定义 1 次微分形式 $\theta$ 为 $\theta=\sum_{j=1}^{n}p_j\,dq_j$，在 $M$ 上的平凡复线丛 $L=M\times\C$ 上引入联络 $\nabla=d-\sqrt{-1}\,\hbar^{-1}\theta$。此时，对 $L$ 的截面即 $M$ 上的复值函数 $s$，令 $\hat f s=\sqrt{-1}\,\hbar\,\nabla_{X_f}s+fs$，把与坐标函数 $q_j,p_j$ 相应确定的 $\hat q_j,\hat p_j$ 的作用写出来，便得到
\[
  \hat p_j=-\sqrt{-1}\,\hbar\,\frac{\partial}{\partial q_j},\qquad
  \hat q_j=\sqrt{-1}\,\hbar\,\frac{\partial}{\partial p_j}+q_j .
\]
因此，特别地，若只看对仅依赖于 $q_1,\dots,q_n$ 的函数的作用，就得到所谓正则量子化的手序。正则量子化本是把 $p_j$ 替换为微分算子 $-\sqrt{-1}\,\hbar\,\frac{\partial}{\partial q_j}$ 的操作。像这样，取 $L$ 的 $L^{2}$ 截面全体的空间作为 Hilbert 空间，就给出量子化。

% [原文 p.xii]
但这个空间作为量子化的空间太大了。下面说明的\textbf{极化}（polarization）给出了抽出恰当空间的方法。

设 $(M,\omega)$ 为 $2n$ 维辛流形。$M$ 的切丛 $TM$ 的复化 $TM_{\C}$ 的可积 Lagrangian 子丛 $V_P\subset TM_{\C}$ 称为极化。这里，子向量丛 $V_P$ 称为\textbf{可积}（integrable）的，是指对 $V_P$ 的截面 $X,Y$，$[X,Y]$ 仍是 $V_P$ 的截面。又，$V_P$ 称为 \textbf{Lagrangian} 的，是指对 $M$ 的任意点 $x$，$V_P$ 在 $x$ 上的纤维维数为 $n$，且在其上 $\omega$ 恒等于 $0$。给定这样的极化时，令
\[
  \mathcal{H}_0=\{s\in\mathcal{H}\mid \nabla_X s=0,\ X\in\Gamma(M,V_P)\},
\]
作为 Hilbert 空间，采用 $\mathcal{H}_0$ 代替 $\mathcal{H}$。上面所述的正则量子化手序中抽出了仅依赖于 $q_1,\dots,q_n$ 的函数，这正由投影 $d\pi\colon TM_{\C}\to TN_{\C}$ 的核所给出的极化的归结。

设 $(M,\omega)$ 为 Kähler 流形。令 $V_P=TM^{0,1}$ 为由 $(0,1)$ 向量全体构成的 $TM_{\C}$ 的子丛，则 $V_P$ 给出上述意义的极化。它称为 Kähler 极化。关于 Kähler 极化的 Hilbert 空间 $\mathcal{H}_0$，无非就是 $L$ 的正则截面全体所成的空间 $H^{0}(M,L)$。作为重要的应用，举出 \textbf{Borel--Weil 理论}。设 $G$ 为紧 Lie 群，$T$ 为其极大环面。$G/T$ 称为旗流形（flag manifold），具有 Kähler 流形的结构。Borel--Weil 理论是说：把 $G$ 的不可约表示实现为 $M$ 上正则复线丛 $L$ 的正则截面全体所成的空间 $H^{0}(M,L)$。这可以看作第 1 章要说明的共形场理论的原型。

\para{sigma 模型}
解说称为 sigma 模型的场论的基本框架。首先固定 Riemann 流形 $M$。取 $X$ 为紧的、带定向的 $d+1$ 维 Riemann 流形，把从 $X$ 到 $M$ 的光滑映射 $\phi$ 在这里称为\textbf{场}（field）。sigma 模型中场的全体的集合，是从 $X$ 到

% [原文 p.xiii]
$M$ 的光滑映射全体所成的空间
\[
  \mathcal{C}_X=\Map(X,M),
\]
即为此。流形 $X$ 是时空的模型，$d+1$ 维这一说法意味着空间 $d$ 维、时间 $1$ 维。一般地，场可以有种种变体。例如在规范理论中，场被看作 $X$ 上纤维丛的联络。对场考虑作用量，理论便由此展开。这里，作用量是 $\mathcal{C}_X$ 上的函数 $S_X\colon\mathcal{C}_X\to\R$，在 sigma 模型中由
\[
  S_X(\phi)=\int_X |d\phi|^{2}\,dv
\]
给出。这里 $dv$ 是表示 $X$ 的 Riemann 度量所对应的体积元素的 $d+1$ 次微分形式。使作用量 $S_X$ 取极小的 $\phi$ 称为调和映射。

下面叙述的 $S_X$ 的几条性质都容易证明，但它们在由此展开的理论中起重要作用。
\begin{enumerate}[label=(\arabic*),leftmargin=2.2em,itemsep=2pt]
\item 设 $f\colon X'\to X$ 为等距变换，则 $f$ 诱导映射 $f^{*}\colon\mathcal{C}_X\to\mathcal{C}_{X'}$，且 $S_{X'}(f^{*}\phi)=S_X(\phi)$ 成立。
\item 把 $X$ 的定向反转所得者记作 $-X$，则 $S_{-X}(\phi)=-S_X(\phi)$，$\phi\in\mathcal{C}_X$ 成立。
\item 设 $X=X_1\sqcup X_2$ 为无公共部分的并集，则对 $\phi\in\mathcal{C}_X$，把到 $X_1,X_2$ 的限制分别记作 $\phi_1,\phi_2$ 时，
\[
  S_X(\phi)=S_{X_1}(\phi_1)+S_{X_2}(\phi_2)
\]
成立。
\item 对 $X$ 的 $d$ 维子流形 $Y$，把 $X$ 分解为 $X=X_+\cup X_-$、$\partial X_+=Y$、$\partial X_-=-Y$ 时，对在 $Y$ 上一致的 $\phi_+\in\mathcal{C}_{X_+}$、$\phi_-\in\mathcal{C}_{X_-}$，存在 $\phi\in\mathcal{C}_X$ 使
\[
  S_X(\phi)=S_{X_+}(\phi_+)+S_{X_-}(\phi_-)
\]
成立。
\end{enumerate}

% [原文 p.xiv]
\para{由 Feynman 路径积分进行的量子化}
在古典力学中，在两点之间运动的粒子的轨道，由使关于 Lagrangian 的作用量积分 $S$ 取极小来确定。与此相对，在量子力学的图景中，对两点之间的所有道路 $\gamma\colon[a,b]\to\R^{n}$，考虑作用量给出的概率函数 $e^{\sqrt{-1}S(\gamma)/\hbar}$ 的贡献；形式上，这用道路空间上的测度 $\mu$ 表示为
\[
  \int e^{\sqrt{-1}S(\gamma)/\hbar}\,d\mu(\gamma).
\]
这类积分称为 \textbf{Feynman 路径积分}（Feynman path integral）。积分的数学严格化，除可归结为 Gauss 积分反复的情形外，是非常困难的。在可把 Planck 常数 $\hbar$ 视为接近 $0$ 的情形下，本质上对路径积分有贡献的是 $S$ 的临界点，即古典力学中粒子的轨道，其余部分则因振动积分相互抵消。这样一来，从量子论便再现出古典论。关于 Feynman 路径积分，也请参考文献 [5]。

取 $X$ 为紧的、无边界、带定向的 $d+1$ 维流形，在 sigma 模型的情形要实行 Feynman 路径积分的量子化，就必须在 $\mathcal{C}_X$ 上引入测度 $\mu_X$，计算配分函数
\[
  Z(X)=\int_{\mathcal{C}_X} e^{\sqrt{-1}S(\phi)/\hbar}\,d\mu_X(\phi).
\]
如上所述，遗憾的是，严格定义测度 $\mu_X$，在现代数学中多数场合尚未做到。这里，我们写出期待配分函数 $Z(X)$ 满足的性质，并以从别的观点构造 $Z(X)$ 为目标。关键的想法是：对带边界的流形定式化理论，从而如上文 (4) 所述把流形分解为两部分，来计算配分函数 $Z(X)$。再继续稍形式化的讨论。设 $X$ 的边界非空，令 $\partial X=Y$。对 $\alpha\in\mathcal{C}_Y$，把到 $Y$ 的限制与 $\alpha$ 一致的 $\phi\in\mathcal{C}_X$ 全体的集合记作 $\mathcal{C}_X(\alpha)$。仍然只是形式的，令

% [原文 p.xv]
\[
  Z(X)_{\alpha}=\int_{\mathcal{C}_X(\alpha)} e^{\sqrt{-1}S(\phi)/\hbar}\,d\mu_X(\phi),
\]
则可知 $Z(X)$ 属于 $C_Y$ 上的函数空间 $Z_Y$。$Z_Y$ 取怎样的函数空间，将在本书处理的 Wess--Zumino--Witten 模型、Chern--Simons 理论中具体考察。说明 $Z_Y$ 可被看作某种复线丛的截面空间，是一个重要目标。$Z_Y$ 称为\textbf{态空间}（quantum Hilbert space）。

由 sigma 模型处所述的性质 (1)--(4)，可以窥知 $Z_Y$ 应满足的性质，它们成为构造理论时的指针。以公理的形式举出其中几条。

\noindent（关于定向的性质）\quad $Z_{-Y}=Z_Y^{*}$。

\noindent 这里 $-Y$ 是给 $Y$ 以相反定向所得者，右端表示对偶空间。

\noindent（乘法性）\quad $Z_{Y_1\sqcup Y_2}=Z_{Y_1}\otimes Z_{Y_2}$。

\noindent（粘合规则）\quad $Z(X)=\bigl(Z(X_+),Z(X_-)\bigr)$。

\noindent 这里 $X=X_+\cup X_-$，$\partial X_+=Y$，$\partial X_-=-Y$，上式右端表示 $Z_Y$ 与 $Z_Y^{*}$ 之间的自然配对。

\para{Wess--Zumino--Witten 模型与环路群}
第 1 章叙述环路群及其上的复线丛，并用它定式化 Wess--Zumino--Witten 模型。设 $\Sigma$ 为紧 Riemann 面，$f\colon\Sigma\to G$ 为光滑映射。这里 $G$ 是紧 Lie 群，本书主要处理 $G=SU(2)$ 的情形。对 $f$ 定义 Wess--Zumino--Witten 作用量 $S_\Sigma(f)$，求其配分函数即在 $f\colon\Sigma\to G$ 的全空间上 $e^{-S_\Sigma(f)}$ 的 Feynman 路径积分，就是理论的目标。这里，对 $\Sigma$ 带边界的情形加以定式化，并说明 $e^{-S_\Sigma(f)}$ 可看作环路群上复线丛的纤维的元素。又，期待配分函数可看作这个复线丛的截面。

环路群上的复线丛，是量子化环路群所得的对象，其截

% [原文 p.xvi]
面空间上，作用着作为环路群无穷小版本的仿射 Lie 代数。目标是从仿射 Lie 代数表示论的立场，规定与配分函数对应的截面。

共形场理论是具有关于共形变换的协变性的、Riemann 面上算子的理论。这里说明可以从 Wess--Zumino--Witten 模型构造这样的共形不变场论。特别地，对 Riemann 球面的情形，描述给出共形块空间维数的融合规则。进一步，说明在由共形块空间构成的、构形空间上的向量丛中，自然地进入称为 KZ 联络的平坦联络。

\para{Jones--Witten 理论}
Jones 多项式是 1980 年代中期由 Jones 从算子代数的观点发现的纽结不变量。本书从作为 KZ 方程的单值出现的辫群表示出发，构造 Jones 多项式。1980 年代后半，Witten 从 Chern--Simons 规范理论的立场提出了三维流形的拓扑不变量，Jones 多项式也随之获得新的解释。设 $M$ 为带定向的三维流形。取 $G$ 为 Lie 群 $SU(2)$，$P$ 为 $M$ 上的 $SU(2)$ 主丛。由于 $SU(2)$ 单连通，$P$ 在拓扑上是 $M$ 上的平凡 $SU(2)$ 丛。把 $P$ 的联络全体的空间记作 $\mathcal{A}_M$。$\mathcal{A}_M$ 成为仿射空间，以平凡联络为基点，把它与取值于 $G$ 的 Lie 代数 $\mathfrak{g}$ 的 $M$ 上 1 次微分形式的空间 $\Omega^{1}(M,\mathfrak{g})$ 等同起来。$P$ 的规范群 $\mathcal{G}$ 可看作从 $M$ 到 $G$ 的光滑映射全体的空间 $\Map(M,G)$。联络的空间 $\mathcal{A}_M$ 上，规范群 $\mathcal{G}$ 通过拉回从右侧作用。

对 $P$ 的联络 $A\in\mathcal{A}_M$，令
\[
  CS(A)=\frac{1}{8\pi^{2}}\int_M \Tr\Bigl(A\wedge dA+\frac{2}{3}A\wedge A\wedge A\Bigr),
  \qquad A\in\mathcal{A}_M,
\]
称 $CS$ 为 \textbf{Chern--Simons 泛函}（Chern--Simons functional）。Witten 给出的三维流形拓扑不变量，用路径积分形式地

% [原文 p.xvii]
表示为
\[
  Z_k(M)=\int_{\mathcal{A}_M/\mathcal{G}} \exp\bigl(2\pi\sqrt{-1}\,k\,CS(A)\bigr)\,DA .
\]
这里 $k$ 是整数。为了把这一表述作为严格的数学来把握，要对 $M$ 带边界的情形定式化理论，并确定相当于态空间的对象。态空间成为共形场理论中的共形块空间，用共形块的单值即可描述粘合规则。这样一来，当给定 Dehn 手术、Heegaard 分解等流形的表示时，就能得到 $Z_k(M)$ 的描述。

把这个方法稍详细地说明，大致如下。对带边界的三维流形 $M$，设 $M$ 的边界为 $\Sigma$，取 $\Sigma$ 上的 $SU(2)$ 联络 $\alpha$。对 $\mathcal{A}_M$ 中限于 $\Sigma$ 时与 $\alpha$ 一致的元素全体，形式地施行上述路径积分，则确定 $Z_k(M)_{\alpha}$ 应属于的态空间。用规范群对 Chern--Simons 泛函的作用来计算，可知态空间可看作 $\Sigma$ 上平坦 $SU(2)$ 向量丛的模空间上复线丛的截面空间。进一步，用 Kähler 极化限制到正则截面的空间，便得到称为共形块空间的复向量空间。这里重要的是：共形块空间成为有限维线性空间。出发点虽是联络全体的空间中的 Feynman 路径积分这一无限维对象，但经过上述手序，就可以抽出有限维线性空间。把 $M$ 分解为以 Riemann 面为边界的两个流形，便可用粘合映射对共形块空间的作用来描述 $Z_k(M)$。这就是上面所述的 $Z_k(M)$ 的 Dehn 手术、Heegaard 分解表示。像这样，从无限维对象中抽出有限的不变量，是本书一贯的主题。

以上是经由共形场理论、严格理解由 Chern--Simons 泛函的路径积分所得不变量的手法；最后一章则举出另一种手法，即 \textbf{Chern--Simons 微扰理论}（Chern--Simons perturbation

% [原文 p.xviii]
theory）。作为有限维的类比，考虑积分
\[
  Z_k=\int_{\R^{n}} e^{\sqrt{-1}k f(x_1,\dots,x_n)}\,dx_1\cdots dx_n
\]
在 $k\to\infty$ 时的渐近行为。在适当假设下，它可分解为来自 $f(x_1,\dots,x_n)$ 的临界点的贡献之和，各项的主项用临界点的 Hessian 表示。进一步，渐近展开的高次项由 Feynman 图描述。把这样的考察应用到 Chern--Simons 泛函的路径积分的情形，用 Feynman 图与 Green 形式构造三维流形的拓扑不变量。

% [Tutor Note: 本导论的叙述顺序（古典力学 → 复线丛与量子化 → sigma 模型 → Feynman 路径积分 → WZW/Chern--Simons）就是全书的概念骨架：第 1 章的共形块、第 2 章的辫群表示与 Witten 不变量、第 3 章的微扰展开都能挂在这条线上。初读可先记住 p.xv 的「态空间 = 复线丛的截面空间」与 p.xvii 的「无限维对象 → 有限维不变量」两句话。]

% [Tutor Note: p.x 中 Dirac 的对应写的是 $\hat f_1\hat f_2-\hat f_2\hat f_1=-\sqrt{-1}\hbar\hat f_3$，而 p.x 末把 $\hat f$ 归一化为 $\sqrt{-1}\hbar\tilde f/\hbar$；两处相差一个常数因子，是作者为把 $\hbar$ 显式提出而作的约定，对易子 $[\hat f_1,\hat f_2]=-\sqrt{-1}\hbar\widehat{\{f_1,f_2\}}$ 与此一致。]

% [Tutor Note: p.xi 的曲率恒等式 $\nabla_X\nabla_Y-\nabla_Y\nabla_X=\nabla_{[X,Y]}-2\pi\sqrt{-1}\omega(X,Y)$ 中，$2\pi\sqrt{-1}$ 的因子来自联络的约定 $\nabla=d-2\pi\sqrt{-1}\alpha$（第一 Chern 形式为 $d\alpha$ 而非 $2\pi$ 倍）。不同教材常把 $2\pi$ 并入 $\omega$，对照时须先确认此约定。]

% [Tutor Note: p.xv 的三条公理只写到「粘合规则」，作者在此省略了「$\dim Z_Y<\infty$」与「$Z(X)$ 对 $X$ 的微分同胚不变」两条；Atiyah 的拓扑场论公理共五条，第 2 章 §2.4 会再用到。]

% [Tutor Note: p.xvi 的 $CS(A)=\frac{1}{8\pi^{2}}\int_M\Tr(A\wedge dA+\frac{2}{3}A\wedge A\wedge A)$ 中，$\frac{2}{3}$ 的系数与 $\frac{1}{8\pi^{2}}$ 的归一化都和 $SU(2)$ 的 Killing 型取法有关；第 2 章 §2.5 与第 3 章 §3.2 沿用同一约定，跨章引用时不要换成别的归一化。]
